\chapter{قبول الصحابة لخبر الآحاد وشروط الحديث المرسل}

\section{رجوع الصحابة وخلفائهم لخبر الآحاد}
الحمد لله والصلاة والسلام على رسول الله، أما بعد؛ فيواصل الإمام الشافعي رحمه الله تعالى إقامة الحجج الدامغة على حجية خبر الواحد (الآحاد) الملزم في العقائد والأحكام، وذلك بسرد وقائع عملية متواترة لكبار الصحابة، رجعوا فيها عن فتاواهم وأقضيتهم واجتهاداتهم بمجرد بلوغهم خبراً لواحد من الصحابة عن رسول الله ﷺ.

\subsection{رجوع عمر بن الخطاب في دية الأصابع والميراث من الدية}
قضى عمر بن الخطاب رضي الله عنه باجتهاده في دية أصابع اليدين تفاضلاً؛ فقضى في الإبهام بخمس عشرة، وفي التي تليها بعشر، وفي الوسطى بعشر، وهكذا (لأن المنافع والجمال تختلف). حتى وُجد كتاب آل عمرو بن حزم الذي كتبه النبي ﷺ لعمرو حين بعثه لليمن، وفيه: «وفي كل إصبع مما هنالك عشر من الإبل»، فرجع عمر رضي الله عنه عن قضائه إلى نص حديث النبي ﷺ. 

وكذلك كان عمر رضي الله عنه يفتي بأن (الدية للعاقلة)، ولا ترث المرأة من دية زوجها المقتول شيئاً، حتى أخبره الضحاك بن سفيان رضي الله عنه أن رسول الله ﷺ كتب إليه: «أن يورّث امرأة أشيم الضبابي من دية زوجها»، فرجع عمر عن رأيه إلى سنة رسول الله ﷺ.

\subsection{رجوعه في دية الجنين وفي أخذ الجزية من المجوس}
ناشد عمر بن الخطاب رضي الله عنه الناس: من سمع من النبي ﷺ في إسقاط الجنين شيئاً؟ (وكاد أن يقضي فيه باجتهاده)، فقام حَمَلُ بن مالك بن النابغة رضي الله عنه فأخبره أنه كان متزوجاً بامرأتين (ضرتين)، فضربت إحداهما الأخرى بعمود خيمة فألقت جنيناً ميتاً، فقضى فيه النبي ﷺ بـ «غُرة» (أي عبد أو أمة). فقال عمر: «لو لم أسمع فيه لقضينا فيه برأينا»، ورجع لحكم النبي ﷺ.
وكذلك توقف عمر في أخذ الجزية من المجوس لأنهم ليسوا أهل كتاب، حتى شهد عبد الرحمن بن عوف رضي الله عنه أنه سمع النبي ﷺ يقول: «سُنّوا بهم سنة أهل الكتاب»، فأخذ منهم الجزية بخبر عبد الرحمن وحده.

\section{توجيه طلب عمر لخبر ثانٍ في بعض الأحاديث}
فإن قيل: قد طلب عمر شاهداً آخر مع أبي موسى الأشعري رضي الله عنه في حديث «الاستئذان ثلاثاً»، فهل هذا لعدم قبوله لخبر الواحد؟ 
يُجاب بأن عمر لم يفعل ذلك رداً لخبر الواحد، وإلا لكان متناقضاً يقبل الآحاد مرات ويرده أخرى! بل فعل ذلك من باب (الاحتياط والتثبت) وسد الباب أمام من قد تسول له نفسه الكذب على رسول الله ﷺ، ولذلك قال لأبي موسى: «أما إني لم أتهمك، ولكن خشيت أن يتقول الناس على رسول الله ﷺ». والاحتياط بطلب التأكيد لا ينفي حجية خبر الواحد القاطع.

\section{أمثلة أخرى لرجوع الصحابة والتابعين للآحاد}
\begin{itemize}
    \item \textbf{رجوع ابن عباس في طواف الوداع للحائض:} كان زيد بن ثابت يفتي بأن الحائض لا تنفر حتى تطوف بالبيت طواف الوداع قياساً على عموم الحجاج، فأحاله ابن عباس إلى امرأة أنصارية سألت النبي ﷺ فأمرها بالنفور وسقوط طواف الوداع عنها، فرجع زيد رضي الله عنه لخبر المرأة الواحدة.
    \item \textbf{رجوع ابن عباس في مسألة الخضر:} كذّب ابن عباس نَوْفاً البَكَالي حين زعم أن موسى صاحب الخضر ليس موسى بني إسرائيل، فلما حدّثه أُبي بن كعب بحديث النبي ﷺ في قصة موسى والخضر، سلم ابن عباس ورجع.
    \item \textbf{رجوع طاووس وابن عمر:} رجع طاووس عن صلاة ركعتين بعد العصر حين نهاه ابن عباس وأخبره بنهي النبي ﷺ عنهما. ورجع ابن عمر عن (المخابرة) وهي كراء الأرض الزراعية لما أخبره رافع بن خديج بنهي النبي ﷺ عنها، وترك ما كان يراه حلالاً ينتفع به لخبر رجل واحد.
    \item \textbf{رجوع القضاة والولاة:} قضى سعد بن إبراهيم بن عبد الرحمن بن عوف في مسألة برأي ربيعة الرأي، فأخبره ابن أبي ذئب بحديث عن النبي ﷺ يخالف قضاءه، فدعا سعد بكتاب القضية فشقه ونقض حكمه، وقال: «أنفذ قضاء رسول الله ﷺ وأرد قضاء سعد».
\end{itemize}
وقد أجمع علماء الأمصار في مكة والمدينة والشام واليمن (كطاووس ووهب بن منبه ومكحول) والكوفة (كابن أبي ليلى والأسود وعلقمة والشعبي) على تثبيت خبر الواحد والعمل به في كل زمان.

\section{شروط العمل بالحديث المرسل عند الشافعي}
الحديث المنقطع، ومنه "المرسل" (وهو ما رفعه التابعي للنبي ﷺ مسقطاً الصحابي)، هل تقوم به الحجة؟ 
أسس الإمام الشافعي قواعد صارمة ومحكمة لقبول المرسل، فبيّن أن مراسيل صغار التابعين لا تُقبل ألبتة لكثرة تساهلهم، أما مراسيل (كبار التابعين) كسعيد بن المسيب، فلا تُقبل أيضاً إلا باجتماع شروط تدل على صحة المخرج، وهي:
\begin{enumerate}
    \item \textbf{أن يسند من وجه آخر:} بأن يُروى نفس الحديث بإسناد متصل (مسند) يوافق معنى المرسل.
    \item \textbf{أن يوافق مرسلاً آخر:} بأن يرسله تابعي آخر كبير أخذ العلم عن شيوخ غير شيوخ التابعي الأول، لانتفاء التواطؤ.
    \item \textbf{موافقة قول صحابي:} أن يوجد من أقوال الصحابة ما يوافق دلالة هذا الحديث المرسل.
    \item \textbf{موافقة عوام أهل العلم:} أن يفتي عوام أهل العلم وعامتهم بمقتضى ما جاء في هذا المرسل.
    \item \textbf{شروط في الراوي المُرسِل نفسه:} أن يكون التابعي المرسِل لا يروي (إذا سمّى) إلا عن الثقات المأمونين، ولا يرغب عن الرواية عن مجهول، وإذا شرك الحُفّاظ المأمونين في حديث لم يخالفهم، بل يوافقهم أو ينقص عنهم ولا يزيد.
\end{enumerate}
ومع كل هذه الشروط، فإن الحديث (المسند المتصل) يظل أقوى من (المرسل)؛ لأن المرسل يحمل انقطاعاً ورواية عن مجهول، والاحتياط للدين يقتضي تقديم المتصل، لكن من توافرت في مرسلاته هذه الشروط من كبار التابعين أُلحقت بالحجة.

\section{خاتمة في خطورة رد السنن}
شنّع الشافعي رحمه الله على من يرد الأحاديث المسندة الصحيحة بالهوى أو العصبية، كمن يتمسك بالأحاديث المرسلة الضعيفة لنصرة مذهبه، ويرد الأحاديث المتصلة الصحيحة التي تلزمه الحجة بها، فهذا من اتباع الهوى والتقليد الأعمى المذموم الذي يودي بصاحبه للضلال.